% Generated draft: requires mathematical and source review.
\documentclass[11pt]{article}
\usepackage[margin=25mm]{geometry}
\usepackage{fontspec}
\setmainfont{Latin Modern Roman}
\usepackage{amsmath,amssymb,mathtools,unicode-math}
\setmathfont{Latin Modern Math}
\usepackage{enumitem,needspace,xurl,hyperref}
\setlength{\parindent}{0pt}
\setlength{\parskip}{5pt}
\setlength{\emergencystretch}{3em}
\newcommand{\N}{\mathbb{N}}
\newcommand{\Z}{\mathbb{Z}}
\newcommand{\Q}{\mathbb{Q}}
\newcommand{\R}{\mathbb{R}}
\newcommand{\C}{\mathbb{C}}
\newcommand{\F}{\mathbb{F}}
\newcommand{\A}{\mathbb{A}}
\newcommand{\PP}{\mathbb{P}}
\newcommand{\E}{\mathbb{E}}
\newcommand{\eps}{\varepsilon}
\newcommand{\dd}{\,\mathrm{d}}
\begin{document}
\section*{UnsolvedMath ID: 30001488}
\textbf{Source problem code:} catalogue entry
\subsection{Definitions and mathematical statement}
A metric space $(X,d)$ is Ptolemy if, for all $a,b,c,e\in X$,
\[
d(a,c)d(b,e)\leq d(a,b)d(c,e)+d(a,e)d(b,c).
\]
Ptolemy equality means equality in this inequality for the specified ordering of the four points. The supplied record does not specify the metric or Möbius structure used on punctured spaces, the definition of an $R$-circle, or the meaning of self-duality of projection. Accordingly, the following records the conjectural target with these qualifications explicit; it is not a complete axiomatization.

Let $X$ be a compact Ptolemy space, with a distinguished family $\mathcal R$ of subsets called $R$-circles, satisfying the four source hypotheses:
\begin{enumerate}
\item Any two distinct points of $X$ belong to some member of $\mathcal R$.
\item A fourth point satisfying Ptolemy equality with three points of an $R$-circle belongs to that circle. The required ordering and distinctness conditions are not supplied.
\item Closest-point projection to an $R$-circle is self-dual. The metric, existence and uniqueness requirements for this projection, and the self-duality identity are not supplied.
\item Every Möbius map between $R$-circles extends to a Möbius automorphism of $X$. The record does not specify whether these maps must be onto the target circles or give a precise definition of the relevant Möbius structures.
\end{enumerate}

For metric spaces $(Y,\rho)$ and $(B,\delta)$, a submetry is a surjective map $\pi:Y\to B$ satisfying
\[
\pi\bigl(\overline B_\rho(y,r)\bigr)=\overline B_\delta(\pi(y),r)
\qquad(y\in Y,\ r\geq0).
\]
Here $\overline B_\rho(y,r)=\{z\in Y:\rho(y,z)\leq r\}$. To apply this definition below, a metric $d_\omega$ on $X\setminus\{\omega\}$ must be specified; the supplied record does not provide it.

\emph{Conjectural target.} For every $\omega\in X$, there should exist a Euclidean space $B_\omega\cong\R^k$ and a submetry
\[
\pi_\omega:(X\setminus\{\omega\},d_\omega)\longrightarrow B_\omega.
\]
Möbius automorphisms should induce compatible homotheties of the bases. The natural equivariance diagram would require, for every Möbius automorphism $g$ of $X$, an induced map $h_{g,\omega}:B_\omega\to B_{g\omega}$ such that
\[
\pi_{g\omega}\circ g=h_{g,\omega}\circ\pi_\omega
\quad\text{on }X\setminus\{\omega\}.
\]
The exact meaning of homothety and the additional compatibility requirements are not supplied, so this equation records the intended form of equivariance without resolving those requirements.

For $b\in B_\omega$, write
\[
F_{\omega,b}=\pi_\omega^{-1}(\{b\}),\qquad
K_{\omega,b}=F_{\omega,b}\cup\{\omega\}.
\]
The completed fibers $K_{\omega,b}$ should be $K$-circles homeomorphic to $S^p$, with
\[
k+p=\dim_\partial X.
\]
The notation $\dim_\partial X$ is not defined in the record, and uniformity of $k$ and $p$ across the choices of $\omega$ is not explicitly specified. Writing $\mathcal K$ for the intended family of $K$-circles, the requested incidence properties are:
\begin{enumerate}
\item[(1K)] For every distinct $x,y\in X$, there is exactly one $K\in\mathcal K$ containing both $x$ and $y$.
\item[(2K)] For every $K\in\mathcal K$ and $R\in\mathcal R$, one has $|K\cap R|\leq2$.
\item[(3K)] For every completed fiber $K_{\omega,b}$ and every $x\in X\setminus K_{\omega,b}$, there is exactly one $R\in\mathcal R$ such that
\[
x,\omega\in R,\qquad R\cap F_{\omega,b}\ne\varnothing.
\]
\item[(4K)] For two distinct fibers $F,F'$ of $\pi_\omega$ and two $R$-circles $R_1,R_2$ through $\omega$ described in the source as meeting $F,F'$, every third fiber $H$ satisfies
\[
H\cap R_1\ne\varnothing\quad\Longleftrightarrow\quad H\cap R_2\ne\varnothing.
\]
The supplied wording does not unambiguously state whether each $R_i$ must meet both $F$ and $F'$, or whether the incidences are assigned separately. This antecedent remains unresolved.
\end{enumerate}
The record also does not explicitly settle whether $\mathcal K$ consists exactly of all completed fibers over all punctures or is a separately specified family represented by them. These missing definitions and incidence qualifications require consultation of the cited report before a fully precise statement can be accepted.
\subsection{Short English statement}
The supplied conjecture asks whether four axioms for $R$-circles in a compact Ptolemy space force Euclidean submetry quotients at every puncture, compatible with Möbius automorphisms, whose completed fibers are spherical $K$-circles satisfying properties (1K)--(4K). Several definitions and the incidence hypothesis in (4K) are missing or unclear in the supplied record. Its open status is supplied dated curation, not an independently verified conclusion.
\subsection{Sources}
\begin{itemize}
\item Local UnsolvedMath registry: supplied record with permanent ID \texttt{30001488}, title \emph{Submetries of Compact Ptolemy Spaces}, exact source problem code \texttt{catalogue entry}, in the \emph{Oberwolfach Reports Open Problems} collection. The supplied proposed year is 2010; no individual proposer is identified.
\item Supplied authoritative source citation: \emph{Geometrie} (2010), Oberwolfach Reports, DOI 10.4171/owr/2010/26, pp.\ 1535--1572. \url{https://doi.org/10.4171/owr/2010/26}. The collection background describes these reports as documenting mathematical workshops and their open-problem discussions. The supplied literature-source label is \emph{OWR 2010/26}.
\item Additional supplied literature source, labeled \emph{Buyalo--Schroeder}: \url{https://arxiv.org/abs/1211.3237}. No title, theorem number, or precise supporting passage is supplied.
\item Supplied literature assessment, checked on \texttt{2026-08-21}: status \emph{open}. The assessment states that existence in property (3K) was still missing in the extracted statement and that related Ptolemy characterizations do not establish this exact $K$-circle submetry assertion. This is supplied dated literature triage derived from the Oberwolfach curation, not independently verified here; time-sensitive status and cited claims require independent verification.
\item The registry labels the statement \texttt{corrected\_verified} and reports that it was checked against the official Oberwolfach report, preserving the stated hypotheses, notation, and open target. This verification claim is supplied metadata and has not been independently confirmed here.
\end{itemize}
\end{document}
